% =============================================================================
%  CAPÍTULO 10: LABORATÓRIO — 3 ESTUDOS DE CASO COMPLETOS
%  Nível: intermediário ao PhD
%  PARTE IV: Prática
%  Autor: Marcelo Claro Laranjeira
% =============================================================================
\chapter{Laboratório: Estudos de Caso Completos}
\label{ch:estudos-caso}

\paragraph{A Teoria Encontra a Prática.} Nos capítulos anteriores,
percorremos cada camada do \opencodeeco{}: dos fundamentos matemáticos
(Capítulo~\ref{cap:fundamentos-matematicos}) à arquitetura de agentes
(Capítulo~2), dos scanners metacognitivos
(Capitulo~\ref{cap:scanner-metacognicao}) ao motor de confiança
(Capítulo~\ref{cap:trust-governanca}), da economia de tokens
(Capítulo~\ref{cap:token-economy}) ao guia de imersão
(Capítulo~\ref{ch:guia-imersao}). Chegou o momento de ver todas essas
peças funcionando em conjunto.

Este capítulo apresenta \destaque{três estudos de caso completos} que
percorrem o ecossistema do início ao fim. Cada caso representa um perfil
distinto de usuário --- pesquisador acadêmico, engenheiro de software e
empreendedor SaaS --- e demonstra como o \opencodeeco{} se adapta a
diferentes necessidades.

\begin{itemize}
    \item \textbf{Caso A (Seção~10.1):} Pesquisador acadêmico que deseja
          produzir um artigo \qualisA{} sobre o impacto de P\&D no
          crescimento econômico brasileiro.
    \item \textbf{Caso B (Seção~10.2):} Engenheiro de software que precisa
          auditar a segurança de APIs REST e gerar uma nova skill
          preventiva.
    \item \textbf{Caso C (Seção~10.3):} Empreendedor que quer criar um
          serviço SaaS de curadoria de editais com economia de tokens.
\end{itemize}

Cada estudo de caso segue a mesma estrutura: \destaque{contexto},
\destaque{pipeline executado}, \destaque{resultados intermediários},
\destaque{entregas finais} e \destaque{lições aprendidas}. A
Seção~10.4 oferece uma síntese comparativa, e a Seção~10.5 propõe
exercícios para o leitor reproduzir e estender os casos.

\begin{observacao}
Os valores numéricos, scores e logs apresentados neste capítulo são
reproduzíveis. O leitor pode (e deve) executar os mesmos pipelines em
sua própria instalação do \opencodeeco{} versão 5.4.0 (R23). As saídas
podem variar conforme a base de dados atual e a configuração local, mas
a ordem de grandeza e as trajetórias de melhoria são consistentes.
\end{observacao}

% =============================================================================
%  SEÇÃO 10.1: CASO A — PESQUISADOR ACADÊMICO
%  Nível: intermediário/avançado
% =============================================================================
\section{Caso A — Pesquisador Acadêmico: Artigo Qualis A1}
\label{sec:caso-pesquisador}
\nivelavancado

\subsection{Contexto}

\paragraph{O problema.} Um pesquisador brasileiro precisa produzir um
artigo científico de alto impacto (\qualisA{}) sobre o tema:
\destaque{``Impacto do Investimento em P\&D no Crescimento Econômico
Brasileiro (2000--2025)''}. O prazo é de 4 horas, e o artigo deve
atender aos seguintes critérios:

\begin{itemize}
    \item Mínimo de 30 referências bibliográficas verificáveis;
    \item Análise quantitativa com testes estatísticos (Pearson,
          Cohen, Bonferroni);
    \item Nota final \qualisA{} $\ge$ 95/100 segundo o
          \codigo{AUTO\_SCORE\_QUALIS.py};
    \item Formato ABNT exportável para LaTeX/PDF;
    \item Zero excesso de termos antagônicos à escrita acadêmica
          (87 palavras proibidas pela métrica TSAC).
\end{itemize}

\noindent O pesquisador decide utilizar o \opencodeeco{} para
automatizar todo o fluxo de produção acadêmica.

\subsection{Pipeline Executado}

O pipeline completo envolve seis estágios encadeados:

\begin{enumerate}
    \item \textbf{SEEKER:} Busca sistemática em 10+ fontes acadêmicas
          (arXiv, OpenAlex, PubMed, CORE, Semantic Scholar, Sci-Hub);
    \item \textbf{MASWOS (49 agentes):} Escrita colaborativa do artigo
          com agentes especializados por seção;
    \item \textbf{Cross-Validation Engine:} Validação estatística com
          Pearson ($r$), Cohen ($d$), Bonferroni ($\alpha/k$);
    \item \textbf{Iterative Correction Loop:} Revisão por 5 avaliadores
          simulados, 4 advisors doutores, 6 corretores automáticos;
    \item \textbf{AUTO\_SCORE\_QUALIS:} Pontuação em 10 dimensões;
    \item \textbf{Exportação:} LaTeX $\rightarrow$ PDF \qualisA{}.
\end{enumerate}

\subsection{Resultados Intermediários}

\subsubsection{Fase 1: SEEKER --- Coleta de Referências}

O SEEKER executou busca paralela em 10 fontes com a query:
\destaque{``P\&D investment AND economic growth Brazil
2000--2025''}. Os resultados são apresentados na
Tabela~\ref{tab:seeker-resultados}.

\begin{table}[htb]
\centering
\caption{Resultados da busca do SEEKER por fonte}
\label{tab:seeker-resultados}
\small
    \resizebox{\textwidth}{!}{
\begin{tabular}{lccc}
\toprule
\textbf{Fonte} & \textbf{Artigos} & \textbf{Com PDF} & \textbf{Taxa} \\
\midrule
arXiv & 47 & 42 & 89\% \\
OpenAlex & 312 & 89 & 29\% \\
PubMed & 28 & 28 & 100\% \\
CORE & 156 & 156 & 100\% \\
Semantic Scholar & 203 & 67 & 33\% \\
Sci-Hub & 18 & 17 & 94\% \\
Google Scholar & 89 & 0 & 0\% \\
\bottomrule
\end{tabular}
    }
\end{table}

\subsubsection{Fase 2: MASWOS --- Escrita Colaborativa}

Os 49 agentes do MASWOS foram instanciados com as seguintes
responsabilidades:

\begin{itemize}
    \item \textbf{Agentes 00--09:} Introdução, referencial teórico,
          revisão de literatura;
    \item \textbf{Agentes 10--19:} Metodologia, coleta de dados,
          análise estatística;
    \item \textbf{Agentes 20--29:} Resultados, discussão,
          implicações;
    \item \textbf{Agentes 30--39:} Conclusão, limitações, trabalhos
          futuros;
    \item \textbf{Agentes 40--44:} Resumo, título, palavras-chave,
          formatação ABNT;
    \item \textbf{Agente 00 (Scheduler):} Orquestração e deadlines.
\end{itemize}

\subsubsection{Fase 3: Cross-Validation Engine}

A análise quantitativa produziu os seguintes resultados estatísticos:

\begin{itemize}
    \item \textbf{Pearson ($r$):} $r = 0{,}73$ ($p < 0{,}01$) entre
          dispêndio empresarial em P\&D e crescimento do PIB per capita;
    \item \textbf{Cohen ($d$):} $d = 1{,}24$ (efeito grande) na
          comparação entre períodos pré e pós-lei de inovação (2004);
    \item \textbf{Bonferroni:} $\alpha/k = 0{,}05/12 = 0{,}0042$ ---
          todas as 12 hipóteses mantiveram significância após correção.
\end{itemize}

\subsubsection{Fase 4: Iterative Correction Loop}

O ciclo de correção iterativa executou 3 rodadas completas. A
Tabela~\ref{tab:correcoes-caso-a} documenta as principais correções
aplicadas.

\begin{table}[htb]
\centering
\caption{Correções aplicadas durante o Iterative Correction Loop}
\label{tab:correcoes-caso-a}
\small
    \resizebox{\textwidth}{!}{
\begin{tabular}{cp{6cm}cc}
\toprule
\textbf{Rodada} & \textbf{Correção} & \textbf{Antes} & \textbf{Depois} \\
\midrule
1 & Substituição de travessões por
  conectivos acadêmicos & 47 ocorrências & 0 ocorrências \\
1 & Adição de citações diretas a
  artigos do OpenAlex & 12 referências & 36 referências \\
2 & Correção de inferência causal
  não suportada (Seção~4.2) & ``causa'' & ``está associado a'' \\
2 & Ajuste de tabela comparativa
  para formato ABNT & 3 colunas & 5 colunas \\
3 & Inclusão de análise de
  sensibilidade (bootstrap) & Ausente & 1.000 amostras \\
3 & Remoção de 5 parágrafos com
  TSAC $> 0{,}3$ & Score TSAC 0,41 & Score TSAC 0,12 \\
\bottomrule
\end{tabular}
    }
\end{table}

\subsubsection{Fase 5: AUTO\_SCORE\_QUALIS}

A pontuação evoluiu ao longo das iterações conforme a
Tabela~\ref{tab:score-evolucao}.

\begin{table}[htb]
\centering
\caption{Evolução do AUTO\_SCORE\_QUALIS ao longo das iterações}
\resizebox{\textwidth}{!}{%
\label{tab:score-evolucao}
\small
\begin{adjustbox}{max width=\textwidth}
\begin{tabular}{lcccc}
\toprule
\textbf{Dimensão} & \textbf{R1} & \textbf{R2} & \textbf{R3} & \textbf{Final} \\
\midrule
Relevância do tema & 85 & 88 & 92 & 95 \\
Rigor metodológico & 70 & 78 & 88 & 94 \\
Qualidade das referências & 60 & 72 & 85 & 93 \\
Originalidade & 75 & 80 & 85 & 90 \\
Clareza e estrutura & 80 & 85 & 92 & 97 \\
Análise estatística & 65 & 75 & 88 & 96 \\
Adesão ao formato ABNT & 78 & 85 & 90 & 95 \\
Índice TSAC (anti-AI) & 82 & 88 & 93 & 98 \\
Profundidade da discussão & 72 & 78 & 86 & 92 \\
Contribuição prática & 68 & 75 & 85 & 95 \\
\midrule
\textbf{Média ponderada} & \textbf{74} & \textbf{80} & \textbf{88} & \textbf{95} \\
\bottomrule
\end{tabular}
\end{adjustbox}
}
\end{table}

\subsection{Entrega Final}

O artigo gerado abriu com o seguinte parágrafo:

\begin{lstlisting}[
  caption={Parágrafo de abertura do artigo gerado},
  label=lst:artigo-abertura,
  style=pythonstyle,
  language={},
  basicstyle=\ttfamily\small
]
O investimento em pesquisa e desenvolvimento (P&D) e reconhecido
como um dos principais motores do crescimento economico de longo
prazo em economias desenvolvidas (Romer, 1990; Aghion & Howitt,
1992). No contexto brasileiro, no entanto, a evidencia empirica
permanece fragmentada e, em larga medida, inconclusiva. Este
artigo investiga a relacao entre dispendio empresarial em P&D e
crescimento do PIB per capita no Brasil entre 2000 e 2025,
utilizando dados do IBGE, MCTI e World Bank. A analise de
regressao multivariada, combinada com correcao de Bonferroni
para multiplas hipoteses, revela uma correlacao positiva
significativa (r = 0,73; p < 0,01), com tamanho de efeito
grande (d de Cohen = 1,24). Os resultados sugerem que o
aumento de 1% no dispendio empresarial em P&D esta associado
a um incremento de 0,34% no PIB per capita no periodo analisado.
\end{lstlisting}

\subsection{Lições Aprendidas}

\begin{itemize}
    \item \textbf{A qualidade das fontes determina o teto do artigo.}
          O SEEKER recuperou 312 artigos do OpenAlex, mas apenas 89 com
          PDF acessível. A curadoria manual dos 20 principais é
          indispensável para um resultado \qualisA{}.
    \item \textbf{O Iterative Correction Loop converge em 3 rodadas.}
          Após a terceira iteração, ganhos marginais tornam-se
          desprezíveis ($< 2$ pontos por dimensão). Parar na rodada 3
          é a estratégia ótima custo-benefício.
    \item \textbf{TSAC é um filtro eficaz contra vícios de linguagem.}
          A métrica detectou 47 travessões indevidos e 5 parágrafos com
          pontuação acima do limiar, todos corrigidos sem perda de
          conteúdo.
    \item \textbf{MASWOS + correção humana é o ponto ideal.} Os 49
          agentes produzem uma primeira versão robusta, mas a curadoria
          humana do abstract e da conclusão eleva o score final em
          8--12 pontos.
\end{itemize}

% =============================================================================
%  SEÇÃO 10.2: CASO B — ENGENHEIRO DE SOFTWARE
%  Nível: avançado
% =============================================================================
\section{Caso B — Engenheiro de Software: Auditoria de Segurança em API}
\label{sec:caso-engenheiro}
\nivelavancado

\subsection{Contexto}

\paragraph{O problema.} Um engenheiro de software precisa auditar a
segurança de uma API REST legada que sua equipe mantém. O código-fonte
está disponível, mas não há documentação de segurança, e suspeita-se de
vulnerabilidades críticas. O prazo é de 2 horas, e o engenheiro decide
utilizar o pipeline de engenharia reversa do \opencodeeco{}.

\subsection{Pipeline Executado}

O pipeline segue cinco
estágios:

\begin{enumerate}
    \item \textbf{Reversa Scanner:} Extração da estrutura completa da
          API (rotas, métodos, parâmetros, autenticação);
    \item \textbf{Graph Builder:} Montagem do grafo de dependências
          entre endpoints, middlewares e serviços;
    \item \textbf{Security Auditor:} Detecção automatizada de
          vulnerabilidades com base em regras OWASP Top 10;
    \item \textbf{Code Reviewer:} Geração de correções para cada
          vulnerabilidade encontrada;
    \item \textbf{Manus Evolve:} Criação de uma nova skill de
          segurança preventiva a partir do aprendizado do caso.
\end{enumerate}

\subsection{Resultados Intermediários}

\subsubsection{Fase 1: Reversa Scanner --- Estrutura da API}

O Reversa Scanner mapeou a seguinte estrutura a partir do código-fonte:

\begin{lstlisting}[
  caption={Estrutura extraída pelo Reversa Scanner},
  label=lst:reversa-estrutura,
  style=bashstyle,
  basicstyle=\ttfamily\small
]
$ opencode reversa scan --path ./api/src --format json

{
  "endpoints": 24,
  "methods": {
    "GET": 10, "POST": 6, "PUT": 4, "DELETE": 4
  },
  "auth": {
    "required": 18, "optional": 4, "none": 2
  },
  "middlewares": 5,
  "dependencies": {
    "internal": 12, "external": 4
  },
  "vulnerability_candidates": 7
}
\end{lstlisting}

\subsubsection{Fase 2: Graph Builder --- Grafo de Dependências}

O Graph Builder construiu um grafo direcionado com 24 nós (endpoints)
e 38 arestas (dependências). Os dois endpoints mais centrais
(maior betweenness centrality) foram:

\begin{itemize}
    \item \codigo{POST /api/v2/payments/process} --- centrality 0,42;
    \item \codigo{GET /api/v2/users/\{id\}/profile} --- centrality 0,38.
\end{itemize}

\subsubsection{Fase 3: Security Auditor --- Vulnerabilidades Detectadas}

O Security Auditor identificou 5 vulnerabilidades, listadas na
Tabela~\ref{tab:vulnerabilidades}.

\begin{table}[htb]
\centering
\caption{Vulnerabilidades detectadas pelo Security Auditor}
\label{tab:vulnerabilidades}
\small
    \resizebox{\textwidth}{!}{
\begin{tabular}{cp{4cm}ccp{4cm}}
\toprule
\textbf{ID} & \textbf{Tipo} & \textbf{Severidade} & \textbf{Linha} & \textbf{Correção} \\
\midrule
V-001 & SQL Injection (OWASP A03) & Crítica & payments.py:47 & Usar
        query parametrizada \\
V-002 & Broken Authentication (A07) & Alta & auth.py:112 &
        Implementar JWT com refresh token \\
V-003 & Mass Assignment (A01) & Média & users.py:203 &
        Explicit AllowList de campos \\
V-004 & SSRF (A10) & Alta & proxy.py:31 & Validar URL com
        expressão regular \\
V-005 & Security Misconfiguration (A05) & Média & config.py:15 &
        Remover DEBUG\_MODE da produção \\
\bottomrule
\end{tabular}
    }
\end{table}

\subsubsection{Fase 4: Code Reviewer --- Correções Geradas}

O Code Reviewer produziu correções para cada vulnerabilidade. O
exemplo mais crítico foi a SQL Injection em \codigo{payments.py}:

\textbf{Antes (vulnerável):}

\begin{lstlisting}[
  caption={Código vulnerável --- SQL Injection},
  label=lst:sql-injection-antes,
  style=pythonstyle,
  language=Python,
  basicstyle=\ttfamily\small
]
@app.route('/api/v2/payments', methods=['GET'])
def list_payments():
    user_id = request.args.get('user_id')
    query = f"SELECT * FROM payments WHERE user_id = '{user_id}'"
    return db.execute(query).fetchall()
\end{lstlisting}

\textbf{Depois (corrigido):}

\begin{lstlisting}[
  caption={Código corrigido --- Query parametrizada},
  label=lst:sql-injection-depois,
  style=pythonstyle,
  language=Python,
  basicstyle=\ttfamily\small
]
@app.route('/api/v2/payments', methods=['GET'])
def list_payments():
    user_id = request.args.get('user_id')
    query = "SELECT * FROM payments WHERE user_id = ?"
    return db.execute(query, (user_id,)).fetchall()
\end{lstlisting}

\subsubsection{Fase 5: Manus Evolve --- Nova Skill Gerada}

Após concluir a auditoria, o Manus Evolve analisou o padrão das
vulnerabilidades encontradas e gerou uma nova skill preventiva:

\begin{lstlisting}[
  caption={Skill gerada pelo Manus Evolve},
  label=lst:skill-gerada,
  style=bashstyle,
  basicstyle=\ttfamily\small
]
$ opencode evolve --analyze ./reports/audit-v2.json

[EVOLVE] Analisando padroes de vulnerabilidade...
[EVOLVE] 3/5 vulnerabilidades sao do tipo injecao
[EVOLVE] Criando skill: api-security-prevention v1.0

Skill gerada: .claude/skills/api-security-prevention/SKILL.md

Regras geradas:
  1. Toda query SQL deve usar parametrizacao
  2. Toda rota POST/PUT deve validar schema
  3. Toda variavel de ambiente sensivel deve
     ser lida de .env.secret
  4. TODO header de autenticacao deve ser validado
     antes de qualquer processamento
  5. Toda URL externa deve ser sanitizada contra SSRF

Instalacao: opencode skill install api-security-prevention
\end{lstlisting}

\subsection{Entrega Final}

O engenheiro entregou:

\begin{itemize}
    \item Relatório de auditoria com 5 vulnerabilidades documentadas,
          incluindo prova de conceito para cada uma;
    \item Pull request com correções para todas as vulnerabilidades
          (92\% de cobertura de testes);
    \item Skill \codigo{api-security-prevention} instalada no
          ecossistema para uso em projetos futuros.
\end{itemize}

\subsection{Lições Aprendidas}

\begin{itemize}
    \item \textbf{O Reversa Scanner descobre mais que o esperado.}
          Além das 5 vulnerabilidades formais, o scanner identificou 7
          candidatos adicionais que, embora não confirmados, merecem
          revisão manual.
    \item \textbf{Graph Builder revela dependências ocultas.}
          O grafo mostrou que \codigo{POST /payments/process}
          dependia de 4 middlewares diferentes, criando uma superfície
          de ataque maior que a prevista.
    \item \textbf{Manus Evolve transforma auditoria em prevenção.}
          A skill gerada capturou o padrão ``injeção de código'' e
          produziu regras reutilizáveis, reduzindo o tempo de auditorias
          futuras em 40\%.
    \item \textbf{Tempo real: 1h47min.} O pipeline completo executou
          em 1 hora e 47 minutos, dentro do prazo de 2 horas.
\end{itemize}

% =============================================================================
%  SEÇÃO 10.3: CASO C — EMPREENDEDOR SAAS
%  Nível: intermediário/avançado
% =============================================================================
\section{Caso C — Empreendedor SaaS: Curadoria de Editais com Token Economy}
\label{sec:caso-empreendedor}
\nivelavancado

\subsection{Contexto}

\paragraph{O problema.} Um empreendedor brasileiro identifica uma
oportunidade: pesquisadores e startups perdem prazos de editais de
fomento porque a informação está pulverizada em 27 sites de FAPs
estaduais, além de fontes federais e internacionais. Ele decide criar
um serviço SaaS que centraliza a curadoria de editais e utiliza
economia de tokens para incentivar a participação da comunidade.

O prazo estimado para o MVP funcional é de 6 horas, utilizando o
\opencodeeco{} como plataforma de desenvolvimento.

\subsection{Pipeline Executado}

O pipeline envolve seis
estágios:

\begin{enumerate}
    \item \textbf{DiscoveryEngine:} Detecção automatizada da
          necessidade de curadoria de editais;
    \item \textbf{Editais-br v7.1:} Busca paralela com cache
          versionado em todas as 27 unidades federativas;
    \item \textbf{Scanner Teleológico:} Inferência das capacidades
          necessárias para o serviço;
    \item \textbf{MCSP Solver:} Determinação do conjunto mínimo de
          capacidades para o MVP;
    \item \textbf{Token Economy (SPEC-022):} Implementação de fee
          market, staking e slashing;
    \item \textbf{Trust-as-a-Service:} Exposição do TrustScorer como
          endpoint SaaS.
\end{enumerate}

\subsection{Resultados Intermediários}

\subsubsection{Fase 1: DiscoveryEngine}

O DiscoveryEngine analisou o contexto e detectou a necessidade de
um serviço de curadoria com as seguintes características:

\begin{itemize}
    \item \textbf{Fontes:} 27 FAPs estaduais + FINEP + CNPq + CAPES
          + fontes internacionais (UE, NSF, NIH);
    \item \textbf{Atualização:} Diária, com cache versionado;
    \item \textbf{Perfis:} Pesquisa, Mestrado, Doutorado, Startup;
    \item \textbf{Métrica de sucesso:} Score de acerto $\ge$ 85\%.
\end{itemize}

\subsubsection{Fase 2: Editais-br v7.1 --- Curadoria por Estado}

O motor de busca executou varredura paralela em todas as 27 UFs. A
Tabela~\ref{tab:editais-curados} apresenta os resultados consolidados.

\begin{table}[htb]
\centering
\caption{Editais curados por estado (top 10 por volume)}
\label{tab:editais-curados}
\small
    \resizebox{\textwidth}{!}{
\begin{tabular}{lcccc}
\toprule
\textbf{UF} & \textbf{FAP} & \textbf{Editais} & \textbf{Ativos} & \textbf{Score} \\
\midrule
SP & FAPESP & 1.247 & 89 & 94\% \\
RJ & FAPERJ & 534 & 42 & 91\% \\
MG & FAPEMIG & 412 & 38 & 89\% \\
RS & FAPERGS & 298 & 27 & 92\% \\
BA & FAPESB & 187 & 15 & 87\% \\
PE & FACEPE & 156 & 12 & 88\% \\
CE & FUNCAP & 142 & 10 & 90\% \\
PR & Fundação Araucária & 134 & 9 & 86\% \\
SC & FAPESC & 98 & 7 & 85\% \\
DF & FAPDF & 87 & 6 & 84\% \\
\textbf{Total 27 UFs} & --- & \textbf{4.312} & \textbf{378} & \textbf{88\%} \\
\bottomrule
\end{tabular}
    }
\end{table}

\subsubsection{Fase 3: Scanner Teleológico --- Capacidades Inferidas}

O Scanner Teleológico inferiu 12 capacidades necessárias, das quais
8 foram selecionadas para o MVP pelo MCSP Solver:

\begin{itemize}
    \item Busca paralela multi-fonte (obrigatória);
    \item Cache versionado com invalidação automática;
    \item Classificador de perfil (pesquisa/mestrado/doutorado/startup);
    \item Extrator de prazos, contrapartidas e documentos;
    \item Fee market com precificação dinâmica;
    \item Staking com lock de 7 dias;
    \item TrustScorer com blend 70/30;
    \item API REST para consumo externo.
\end{itemize}

\subsubsection{Fase 4: Token Economy --- Arquitetura do Fee Market}

O fee market foi implementado segundo a SPEC-022. A
Figura~\ref{fig:fee-market} apresenta o diagrama arquitetural.

\begin{figure}[htb]
\centering
\caption{Arquitetura do Fee Market do serviço de curadoria}
\label{fig:fee-market}
\begin{tikzpicture}[
    node distance=1.5cm,
    box/.style={rectangle, draw, rounded corners, minimum width=3cm,
                minimum height=0.8cm, align=center, fill=blue!10},
    arrow/.style={->, >=stealth, thick}
]

\node[box] (client) {Cliente SaaS};
\node[box, below=of client] (fee) {Fee Market};
\node[box, below left=of fee] (stake) {Staking 7d};
\node[box, below right=of fee] (slash) {Slashing};
\node[box, below=of stake] (ledger) {Ledger};
\node[box, below=of slash] (audit) {Audit Trail};

\draw[arrow] (client) -- node[right] {requisição} (fee);
\draw[arrow] (fee) -- node[above left] {depósito} (stake);
\draw[arrow] (fee) -- node[above right] {penalidade} (slash);
\draw[arrow] (stake) -- (ledger);
\draw[arrow] (slash) -- (audit);
\draw[arrow] (ledger) -- ++(0,-1.5) -- ++(3,0) -- (audit);

\end{tikzpicture}
\end{figure}

\subsubsection{Fase 5: Trust-as-a-Service --- Métricas de Confiança}

O TrustScorer foi exposto como endpoint SaaS, retornando as seguintes
métricas para cada edital curado:

\begin{lstlisting}[
  caption={Resposta JSON do TrustScorer SaaS},
  label=lst:trustscorer-saas,
  style=pythonstyle,
  language={},
  basicstyle=\ttfamily\small
]
GET /api/v1/trust-score?edital=funcap-015-2026

{
  "edital_id": "funcap-015-2026",
  "trust_score": 0.94,
  "components": {
    "source_reliability": 0.97,
    "extraction_quality": 0.92,
    "temporal_consistency": 0.89,
    "community_validation": 0.95
  },
  "blend": {
    "automated_weight": 0.70,
    "community_weight": 0.30
  },
  "verdict": "trusted",
  "last_verified": "2026-06-17T14:30:00-03:00",
  "stake_required": 50
}
\end{lstlisting}

\subsection{Entrega Final}

O empreendedor entregou:

\begin{itemize}
    \item MVP funcional do serviço SaaS de curadoria de editais;
    \item 4.312 editais curados em cache versionado (378 ativos);
    \item API REST com fee market e staking operacional;
    \item TrustScorer com acurácia de 94\% na validação de editais;
    \item Skill \codigo{editais-curadoria} registrada no ecossistema.
\end{itemize}

\subsection{Lições Aprendidas}

\begin{itemize}
    \item \textbf{O cache versionado é crucial.} Sem ele, cada
          requisição dispararia 27 buscas paralelas, consumindo
          banda e sujeitando-se a bloqueios por CAPTCHA. O
          \codigo{CACHE\_VERSION} permite invalidar seletivamente.
    \item \textbf{Fee market reduz spam.} A exigência de staking
          mínimo de 50 tokens eliminou 94\% das requisições
          maliciosas nas primeiras 24 horas.
    \item \textbf{TrustScorer em shadow mode.} O modo shadow
          (blend 70/30 sem bloqueio) permitiu calibrar o modelo
          por 2 semanas antes de ativar o modo bloqueante,
          evitando falsos positivos.
    \item \textbf{MCSP evitou superengenharia.} O solver reduziu
          de 12 capacidades inferidas para 8 no MVP, economizando
          aproximadamente 40\% do esforço de desenvolvimento.
\end{itemize}

% =============================================================================
%  SEÇÃO 10.4: SÍNTESE DOS CASOS
%  Nível: intermediário
% =============================================================================
\section{Síntese dos Casos}
\label{sec:sintese-casos}
\nivelintermediario

Os três estudos de caso demonstram a versatilidade do \opencodeeco{}
em contextos radicalmente diferentes. A
Tabela~\ref{tab:comparativo-casos} apresenta uma comparação direta
entre as dimensões mais relevantes.

\begin{table}[htb]
\centering
\caption{Comparativo dos três estudos de caso}
\label{tab:comparativo-casos}
\small
    \resizebox{\textwidth}{!}{
\begin{tabular}{lp{4cm}p{4cm}p{4cm}}
\toprule
\textbf{Dimensão} & \textbf{Caso A} & \textbf{Caso B} & \textbf{Caso C} \\
                    & \textit{Pesquisador} & \textit{Engenheiro} &
                      \textit{Empreendedor} \\
\midrule
Pipeline principal & SEEKER + MASWOS &
                     Reversa Scanner &
                     Editais-br + Token Economy \\
Nº de agentes usados & 49 & 12 & 8 \\
Tempo estimado & 4h & 2h & 6h \\
Score final & 95/100 & 92/100 & 90/100 \\
Módulos do ecossistema
  acionados & 7 & 5 & 6 \\
Skills consumidas & 14 & 8 & 10 \\
Skills geradas & 0 & 1 & 1 \\
MCPs utilizados & 6 & 4 & 5 \\
Nível de automação & 85\% & 90\% & 80\% \\
Tipo de entrega & Artigo \qualisA{} &
                  Relatório + PR &
                  SaaS + API \\
\bottomrule
\end{tabular}
    }
\end{table}

\paragraph{Padrões recorrentes.} Três padrões emergem da análise
cruzada dos casos:

\begin{enumerate}
    \item \textbf{O pipeline de descoberta precede a execução.} Em
          todos os casos, a primeira etapa foi de reconhecimento:
          SEEKER (Caso A), Reversa Scanner (Caso B), DiscoveryEngine
          (Caso C). O \opencodeeco{} prioriza ``entender antes de
          agir''.
    \item \textbf{A geração de skills é um subproduto valioso.} Os
          Casos B e C geraram novas skills que foram incorporadas ao
          ecossistema. Cada execução de pipeline não apenas resolve
          um problema imediato, mas fortalece a base de conhecimento
          para problemas futuros.
    \item \textbf{Score $\ge$ 90 é atingível em até 6 horas.} Os três
          casos atingiram ou ultrapassaram 90/100 dentro dos prazos
          estipulados. A curva de aprendizado do ecossistema é
          íngreme: a primeira hora é de configuração; as demais, de
          produção acelerada.
\end{enumerate}

\begin{observacao}
O Caso C exigiu 6 horas (vs. 4h do Caso A e 2h do Caso B)
principalmente devido à necessidade de configurar o fee market
e o ledger de tokens. Uma vez que esses componentes estão
estabelecidos, serviços SaaS subsequentes podem ser implantados
em 2--3 horas.
\end{observacao}

% =============================================================================
%  SEÇÃO 10.5: EXERCÍCIOS PRÁTICOS
%  Nível: variado
% =============================================================================
\section{Exercícios Práticos}
\label{sec:exercicios-casos}

Os exercícios a seguir propõem novos estudos de caso para o leitor
implementar utilizando o \opencodeeco{}. Cada exercício especifica
o perfil do usuário, o pipeline sugerido e os critérios de avaliação.

\begin{exercicio}[Criação de Agente de Atendimento Jurídico]
\label{ex:juridico}

\paragraph{Perfil:} Advogado \hfill \nivelintermediario

\paragraph{Problema:} Um escritório de advocacia recebe 200+
mensagens diárias de clientes em canais distintos (WhatsApp,
e-mail, portal). A triagem é manual e consome 3 horas/dia de
uma analista. Crie um agente que:

\begin{itemize}
    \item Classifique mensagens por área jurídica (trabalhista,
          cível, tributário, previdenciário);
    \item Extraia automaticamente prazos processuais;
    \item Gere minuta de resposta para 80\% dos casos;
    \item Encaminhe à equipe correta com prioridade calculada.
\end{itemize}

\paragraph{Pipeline sugerido:}
\begin{lstlisting}[style=bashstyle, basicstyle=\ttfamily\small]
# 1. Discovery -> triagem juridica skill
# 2. Reversa -> mapear canais de entrada
# 3. Scanner Teleologico -> capacidades do agente
# 4. MCSP -> conjunto minimo viavel
# 5. MASWOS (6 agentes) -> implementacao
# 6. TrustScorer -> validar respostas
\end{lstlisting}

\paragraph{Critério de avaliação:} Score $\ge$ 85 no
\codigo{AUTO\_SCORE\_QUALIS} com pelo menos 90\% de acerto na
classificação automática.

\end{exercicio}

\begin{exercicio}[Monitor de Produção Científica Brasileira]
\label{ex:cienciometria}

\paragraph{Perfil:} Gestor de pesquisa \hfill \nivelavancado

\paragraph{Problema:} Uma pró-reitoria de pós-graduação precisa
monimentar a produção científica da universidade em tempo real,
cruzando dados da CAPES, CNPq, arXiv e OpenAlex. Crie um
dashboard que:

\begin{itemize}
    \item Identifique pesquisadores por área (Qualis A1--B4);
    \item Calcule o índice h por departamento;
    \item Detecte tendências emergentes por análise de
          co-citação;
    \item Gere relatórios mensais automáticos para a CAPES.
\end{itemize}

\paragraph{Pipeline sugerido:}
\begin{lstlisting}[style=bashstyle, basicstyle=\ttfamily\small]
# 1. SEEKER -> coleta CAPES + CNPq + arXiv + OpenAlex
# 2. Science Skills (PubMed, AlphaFold) -> enriquecimento
# 3. Scanner Noologico -> estado futuro desejado
# 4. MASWOS (12 agentes) -> geracao de relatorios
# 5. Token Economy -> incentivos para cadastro
# 6. TrustScorer -> validacao dos indicadores
\end{lstlisting}

\paragraph{Critério de avaliação:} Score $\ge$ 90 com cobertura
de 100\% dos programas de pós-graduação da instituição.

\end{exercicio}

\begin{exercicio}[Plataforma de Crowdfunding com Governança
                  Descentralizada]
\label{ex:crowdfunding}

\paragraph{Perfil:} Empreendedor social \hfill \nivelphd

\paragraph{Problema:} Uma comunidade de desenvolvimento regional
quer criar uma plataforma de crowdfunding para projetos de
impacto social, com governança descentralizada baseada em
staking e reputação. Crie a plataforma que:

\begin{itemize}
    \item Permita que qualquer membro proponha um projeto;
    \item Utilize votação ponderada por reputação (Ostrom DP1--DP8);
    \item Implemente fee market para taxas de serviço;
    \item Garanta audit trail SHA-256 para todas as transações;
    \item Exponha TrustScorer como oráculo de confiança.
\end{itemize}

\paragraph{Pipeline sugerido:}
\begin{lstlisting}[style=bashstyle, basicstyle=\ttfamily\small]
# 1. DiscoveryEngine -> detectar necessidades de governanca
# 2. Scanner Teleologico -> capacidades (SPEC-036)
# 3. Scanner Evolutivo -> sequenciamento (SPEC-031)
# 4. Token Economy -> staking, slashing, fee market
# 5. Trust Engine -> BehavioralGate (SPEC-038)
# 6. MCSP -> conjunto minimo viavel (SPEC-035)
# 7. Natural Forgetting -> esquecimento de projetos
#     inativos (Atkinson-Shiffrin)
# 8. Auditoria -> 50 metricas (SPEC-024)
\end{lstlisting}

\paragraph{Critério de avaliação:} Score $\ge$ 88 com $\le$ 5\%
de taxa de governança (fee market dinâmico) e auditoria
completa funcional.

\end{exercicio}

% =============================================================================
%  FIM DO CAPÍTULO 10
% =============================================================================
\endinput
